% !PRINTED 643
\weberpage{643}

%<<CONT
\gesperrt{is a rational function of degree $n$ of $t$, whose coefficients $\overline{A}_1, \overline{A}_2, \ldots, \overline{A}_n$ are functionals of the field $\overline{Z}$.}

We define the norm and the trace of the functional as in § 172:
\weq{7}{N(\overline{\omega}) = \left|\begin{array}{llll}
\overline{x}_{1,1}, & \overline{x}_{1,2}, & \ldots, & \overline{x}_{1,n} \\
\overline{x}_{2,1}, & x_{2,2}, & \ldots, & \overline{x}_{2,n} \\
\cdot & \cdot & \cdots & \cdot \\
\overline{x}_{n,1}, & \overline{x}_{n,2}, & \ldots, & \overline{x}_{n,n}
\end{array}\right|} % [sic] second entry of second row printed without overline, reproduced as scanned
\weq{8}{S(\overline{\omega}) = \overline{x}_{1,1} + \overline{x}_{2,2} + \cdots + \overline{x}_{n,n},}
and the theorems § 172, (6) and (21) hold unchanged for the norms and traces of the functionals as well.

Then, as in § 172, (8),
\weq{9}{\varphi(t) = N(t - \overline{\omega}),}
and the \gesperrt{discriminant of a system of functionals} is defined as in § 173, (1); and if the functional $\overline{\omega}$ is contained in a field $\overline{\Omega}_1$ of degree $e$ that is contained in $\overline{\Omega}$, then $\overline{\omega}$ satisfies an equation of degree $e$, $\varphi_1(t) = 0$. As in § 172, (19) it is proved that
\[
\varphi(t) = [\varphi_1(t)]^f \quad (n = ef)
\]
is a power of $\varphi_1(t)$.

\begin{center}
§ 179.\\
\gesperrt{Integral Functionals of the Field $\overline{\Omega}$.}
\end{center}

6. \gesperrt{Definition. A functional $\overline{\omega}$ in $\overline{\Omega}$ is called integral if it satisfies an equation:}
\weq{1}{F(\overline{\omega}) = \overline{\omega}^m + \overline{A}_1\overline{\omega}^{m-1} + \cdots + \overline{A}_{m-1}\overline{\omega} + \overline{A}_m = 0}
\gesperrt{whose coefficients $\overline{A}_1, \ldots, \overline{A}_m$ are integral rational functionals.}

From \gesperrt{Gauss's} theorem (Bd. I, § 2) it follows, since by § 177 one knows the prime factors in $\overline{Z}$, that if $F(t)$ is reducible in $\overline{Z}$, every function dividing $F(t)$ also has integral coefficients, and in particular therefore also the equation of lowest degree $\varphi(\overline{\omega}) = 0$ which $\overline{\omega}$ satisfies.

The necessary and sufficient condition for a functional $\overline{\omega}$ to be integral is therefore that
\[
N(t - \overline{\omega})
\]
when ordered according to powers of $t$, has integral rational functionals for coefficients.
